% =============================================================================
% ISI Country Brief — Slovakia
% External Supplier Concentration Profile
%
% ISI Paper Series — Country Brief
% International Sovereignty Institute — Research Unit
% March 2026
% =============================================================================
\documentclass[12pt,a4paper]{article}

% --- ISI Paper Series shared template ----------------------------------------
\usepackage{isi-series}

% --- PDF metadata ------------------------------------------------------------
\hypersetup{
  pdftitle={External Supplier Concentration Profile: Slovakia — ISI EU-27 2024},
  pdfauthor={International Sovereignty Institute -- Research Unit},
  pdfcreator={International Sovereignty Institute},
  pdfsubject={ISI Paper Series -- Country Brief SK, EU-27, Vintage 2024},
  pdfkeywords={sovereignty index, EU-27, supplier concentration,
    slovakia, country brief},
}

% ==============================================================================
\begin{document}
\hypersetup{pageanchor=false}

% ---------- Title Page with Metadata ------------------------------------------
\begin{titlepage}
\centering
\vspace*{2cm}
{\Large\textsc{ISI Paper Series --- Country Brief}}\\[1.2cm]
{\LARGE\bfseries External Supplier Concentration Profile:\\[0.3em]
Slovakia}\\[1.5cm]
{\large International Sovereignty Index (ISI)\\[0.3em]
EU-27 --- Vintage 2024 --- Methodology v1.0}\\[1.2cm]
{\large International Sovereignty Institute --- Research Unit}\\[0.8cm]
{\large March 2026}

\vspace{0.8cm}
\noindent\rule{\textwidth}{0.4pt}
\vspace{0.4cm}

% --- Research Metadata --------------------------------------------------------
\begin{flushleft}
\noindent\textbf{Data basis:} ISI Paper~2 --- EU-27 Empirical Results 2024 (Methodology~v1.0, Vintage~2024).\\[4pt]
\noindent\textbf{Methodology:} ISI Paper~1 --- Methodological Foundations, Version~1.0.

\medskip

\noindent\textbf{JEL Codes:} F02, F15, F52, O30, Q43\\[4pt]
\noindent\textbf{Keywords:} sovereignty index; EU-27; supplier concentration; composite indicator; variance decomposition; defence dependence

\medskip

\noindent\textbf{Related publications in the ISI Paper Series:}
\begin{itemize}[nosep,leftmargin=1.5em]
\item ISI Paper~1 --- Technical Specification (v1.0) --- doi:\\
  \url{https://doi.org/10.5281/zenodo.18764227}
\item ISI Paper~2 --- EU-27 Empirical Results 2024 (v1.0) --- doi:\\
  \url{https://doi.org/10.5281/zenodo.18764170}
\end{itemize}
\end{flushleft}
\vfill
\end{titlepage}

\hypersetup{pageanchor=true}
\pagenumbering{arabic}
\pagestyle{fancy}
\fancyhead[L]{\small\sffamily\textit{ISI Country Brief}}
\fancyhead[R]{\small\sffamily\textit{Slovakia}}
\renewcommand{\headrulewidth}{0.4pt}

% ---------- Body --------------------------------------------------------------

\section{Executive Summary}
\label{sec:summary}

Slovakia ranks~26th in the EU-27 with a composite \ISI score of~0.236, classified as \emph{mildly concentrated}.  The dominant axes are logistics~(0.486) and energy~(0.400).  The structural profile is characterised as dual-axis concentration (logistics and energy).  Under Dirichlet weight perturbation (10\,000~Monte Carlo draws), Slovakia's mean rank is~25.74 with a standard deviation of~1.58; the 5th--95th percentile band spans ranks~22 to~27.


\section{Country Position in the EU-27}
\label{sec:position}

Slovakia occupies rank~26 of~27 EU member states (composite score~0.236).  The EU-27 mean composite score is~0.344 with a standard deviation of~0.070.  Slovakia's score lies below the EU-27 mean, indicating below-average external supplier concentration.  In the ranking, Slovakia is situated between Latvia (rank~25; 0.247) and France (rank~27; 0.236).


\section{Axis Profile}
\label{sec:axis-profile}

\Cref{tab:axis-profile} presents Slovakia's scores on all six \ISI axes.

\begin{table}[htbp]
\centering\small
\caption{Slovakia's \ISI axis profile.}
\label{tab:axis-profile}
\begin{tabular}{@{}lr@{}}
\toprule
\textbf{Axis} & \textbf{Score} \\
\midrule
    Finance           & 0.162 \\
    Energy            & 0.400 \\
    Technology        & 0.217 \\
    Defence           & 0.000 \\
    Critical Inputs   & 0.155 \\
    Logistics         & 0.486 \\
\bottomrule
\end{tabular}
\end{table}

The highest axis score is logistics~(0.486), which lies below the EU-27 axis mean of~0.518.  The second-highest axis is energy~(0.400), above the EU-27 axis mean of~0.365.  The lowest axis score is defence~(0.000), below the EU-27 mean of~0.573.  \emph{Note:}  Slovakia's defence axis score of~0.000 is a data-window artefact --- no qualifying bilateral defence-trade flow exceeded the reporting threshold during the reference period (2019--2023).  This zero does not imply the absence of defence imports; it reflects the structure of the available data.


\section{Structural Drivers}
\label{sec:drivers}

\begin{sloppypar}
Slovakia's composite score is driven by two axes in combination: logistics~(0.486) and energy~(0.400).  The gap between these two axes is~0.086, indicating a dual-axis concentration profile.  Neither axis alone dominates the composite; rather, the combination of elevated scores on both dimensions determines Slovakia's position in the ranking.
\end{sloppypar}


\section{Comparative Context}
\label{sec:comparative}

Across the EU-27, the covariance-based variance decomposition shows that defence (35.1\,\%) and logistics (34.3\,\%) together account for 69.4\,\% of total cross-sectional composite variance.  Slovakia's defence score~(0.000) lies below the EU-27 defence mean~(0.573), and its logistics score~(0.486) lies below the EU-27 logistics mean~(0.518).  Under Dirichlet weight perturbation, Slovakia's rank standard deviation is~1.58, indicating moderate rank stability.


Slovakia's composite of~0.236 is shaped by a structural artefact on the defence axis: the zero score reflects the absence of qualifying bilateral defence-trade flows above the reporting threshold during the 2019--2023 reference period, rather than a literal absence of defence imports.  Because defence accounts for 35.1\,\% of EU-wide composite variance, a zero on this axis exerts a powerful downward pull on the composite.

The remaining five axes are not uniformly low.  Logistics~(0.486) and energy~(0.400) both score near or above their EU-27 means, and technology~(0.217) exceeds the EU-27 technology mean.  Were the defence axis excluded or imputed at the EU-27 mean, Slovakia's composite would shift materially upward---a counterfactual that illustrates the sensitivity of the index to data-window effects on individual axes.

The rank standard deviation of~1.58 and the five-position percentile band~(22--27) are wide relative to countries at similar composite levels.  This volatility arises because weight perturbations that reduce the defence share allow Slovakia's above-mean scores on other axes to gain influence, pulling the rank upward.  Conversely, perturbations that increase the defence share lock the rank at or near the bottom of the distribution.


\section{Interpretation Boundary}
\label{sec:interpretation}

The \ISI measures concentration of external suppliers, not sovereignty in a normative or political sense.  High concentration may reflect geography, economic specialisation, or alliance structures.  A high composite score does not imply policy failure; a low score does not imply policy success.  The index provides an empirical measurement basis --- what policymakers derive from it is a separate question.


% ---------- References --------------------------------------------------------
\clearpage
\vspace*{1.5cm}

\begingroup\emergencystretch=2em\hbadness=10000
\section*{References}

\begin{enumerate}[label={[\arabic*]},leftmargin=2em,itemsep=6pt]
\item International Sovereignty Institute.
  \textit{International Sovereignty Index (ISI) technical specification, version~1.0}.
  Zenodo, 2026.\\ISI Paper Series, Paper~1.
  doi:\,\href{https://doi.org/10.5281/zenodo.18764227}{10.5281/zenodo.18764227}.

\item International Sovereignty Institute.
  \textit{International Sovereignty Index (ISI): EU-27 empirical results 2024 (v1.0)}.
  Zenodo, 2026.\\ISI Paper Series, Paper~2.
  doi:\,\href{https://doi.org/10.5281/zenodo.18764170}{10.5281/zenodo.18764170}.
\end{enumerate}
\endgroup

\end{document}
